\documentclass[11pt]{article}
\usepackage[a4paper,margin=22mm]{geometry}
\usepackage{amsmath,amssymb,amsthm,mathtools,bm}
\usepackage{booktabs,array,microtype,xurl,hyperref}
\hypersetup{colorlinks=true,linkcolor=blue,citecolor=blue,urlcolor=blue,
 pdftitle={P54 fixed-VEV full complex doublet matching}}
\setlength{\emergencystretch}{2em}
\setcounter{tocdepth}{1}
\newcommand{\tr}{\operatorname{Tr}}
\newcommand{\MS}{\overline{\mathrm{MS}}}
\newcommand{\norm}[1]{\left\lVert#1\right\rVert_2}
\newtheorem{proposition}{Proposition}
\title{Route-F P54: Fixed-VEV Tadpoles,\\the Full Complex Doublet Matrix,\\and the Pati--Salam Matching Interface}
\author{P54PQ-v2 four-dimensional mainline}
\date{5 September 2026}
\begin{document}
\maketitle
\begin{abstract}
We fix a symmetry-preserving finite tadpole prescription and compute every
entry of the bosonic $4\times4$ complex doublet curvature, retaining the
real and imaginary copy mixing. At the historical tree background the
retuned matrix has one zero eigenvalue and three positive heavy
eigenvalues. Its light vector rotates by $0.0329814$ radians. The coarse
heavy-gap norm test fails, but an exact relative-matrix certificate proves
positivity at the stated truncation. Neither higher-loop control nor
all-sector quantum stability follows.

The subsequent matching audit finds an independent error in the historical
P2 interface: its PS beta table and threshold projectors do not describe
the same active spectrum. The $SU(2)_L$ matching logarithm must have slope
$-71$, whereas the historical implementation has zero slope. We derive a
same-field-content repair of all six logarithmic identities. Finite
matching, the complete P54 Yukawa evolution and a physical constrained
global fit remain uncomputed; no fit or model exclusion is fabricated.
\end{abstract}
\tableofcontents
\clearpage

\section{Frozen action, background and scope}
The action remains P54PQ-v2, with fields
$3\,16_F+54_{H,\mathbb R}+126_{H,\mathbb C}+10_{H,\mathbb C}+1_{H,\mathbb C}$.
We denote them by $\Psi,\Phi,\Sigma,\phi,S$. Their PQ charges are
$(-1,0,+2,-2,-4)$, and the two symmetric family matrices enter as
\begin{equation}
-\mathcal L_Y=\tfrac12 (Y_{10})_{ij}(\Psi_i\Psi_j)_{10}\phi^*
+\tfrac12 (Y_{126})_{ij}(\Psi_i\Psi_j)_{\overline{126}}\Sigma+\mathrm{h.c.}
\label{eq:Yukawa}
\end{equation}
Scalar conjugation in these two terms differs. The scalar implementation
uses 328 canonical real coordinates, with complex coefficients
$(x_R+i x_I)/\sqrt2$. The chosen background is
\begin{align}
\Phi_0&=\omega\,\mathrm{diag}(-\tfrac25 I_6,\tfrac35 I_4),\\
\Sigma_0&=\sigma\,\Omega_{126},& S_0&=v_s/\sqrt2, &\phi_0&=0,\\
\sigma/\omega&=0.1265,&v_s/\omega&=0.25,&g_U&=0.562602889985.
\end{align}
We use $\omega=1$ for numerical tensors and $\mu=g_U\omega$ for the
\emph{value} of the subtraction scale. All field derivatives hold $\mu$
and renormalized couplings fixed. The historical scale
$M_U=8.5347508424\times10^{14}\,\mathrm{GeV}$ is not recertified here:
Section~\ref{sec:PS} finds a defect in its matching calculation.

The bosonic hard functional is the previously declared background-field
Landau-$\MS$ functional,
\begin{equation}
 V_{1,B}^{\rm hard}=\frac1{64\pi^2}
 \left\{\tr_{290}F_{3/2}(H(x))+3\tr_{33}F_{5/6}(M_V^2(x))\right\},
 \quad F_c(t)=t^2[\log(t/\mu^2)-c].
\label{eq:CW}
\end{equation}
Here $H=V_0''$ is the scalar operator and
$(M_V^2)_{AB}=g_U^2(T_Ax)^T(T_Bx)$. The 38 soft scalar and 12 soft vector
modes are excluded by a gapped spectral branch, not by differentiating a
fixed diagonal list. The vector coefficient includes the gauge/ghost
prescription of this hard potential. This is a zero-momentum matching
curvature, not a pole mass and not a gauge-independent observable.

\section{Step 1: an invariant fixed-VEV tadpole prescription}
\subsection{Why exactly three radial conditions suffice at this vacuum}
Let $R=(\partial_\omega x_0,\partial_\sigma x_0,\partial_{v_s}x_0)$.
The canonical metric gives
\begin{equation}
 R^TR=\mathrm{diag}(12/5,2,1).
\end{equation}
The SM-singlet tangent has five real dimensions: one in $54$, a complex
singlet in $126$, and the complex $S$. Its two phase directions are the
broken abelian gauge orbit and the PQ orbit. They span the two phase
tangents because the gauge orbit acts on $\Sigma_0$ and PQ also acts
nontrivially on $S_0$. There is no additional physical singlet phase at
this background. An SM-invariant gradient has no component in nontrivial
SM representations; invariance under the two phase orbits removes the
other two singlet components. Thus only the three radial tadpoles remain.
This argument uses a symmetry-preserving regulator and the perturbative
potential; it does not erase the nonperturbative PQ anomaly.

Write $t_\alpha=\partial_\alpha V_1^{\rm hard}(x_0)$. Choose the finite
counterterm within the already present invariant quadratic operators:
\begin{equation}
 \delta V=-\frac{\delta\mu^2}{2}\Phi_{ij}\Phi_{ij}
 -\frac{\delta\nu^2}{2\cdot5!}\Sigma_{ijklm}\Sigma^*_{ijklm}
 -\delta\mu_s^2 S^*S.
\end{equation}
Restricted to the vacuum,
\begin{equation}
 \delta V(x_0)=-\frac65\delta\mu^2\omega^2
 -\frac12\delta\nu^2\sigma^2-\frac12\delta\mu_s^2 v_s^2.
\end{equation}
Solving $\partial_\alpha(V_1+\delta V)=0$ gives
\begin{equation}
 \boxed{\delta\mu^2=\frac{5t_\omega}{12\omega},\qquad
 \delta\nu^2=\frac{t_\sigma}{\sigma},\qquad
 \delta\mu_s^2=\frac{t_s}{v_s}.}
\label{eq:tadpole}
\end{equation}
This defines the finite mass-renormalization convention at the chosen
scale. It must be used also in subsequent thresholds; it is not an
additional change in field content.

The canonical real Hessian is
\begin{equation}
 \boxed{H_{\rm ct}=-\delta\mu^2P_{54}
 -\tfrac12\delta\nu^2P_{126}-\delta\mu_s^2P_S.}
\label{eq:CT}
\end{equation}
The factor $1/2$ for $126$ follows from the declared kinetic/tensor
normalization. Replacing it by one fails the phase Ward identity.
For an antisymmetric generator $T$, invariance implies
$\nabla V(x)^TTx=0$. Differentiation gives
\begin{equation}
 H(x)Tx=T\nabla V(x).
\end{equation}
Consequently the loop curvature along a broken orbit need not vanish
before tadpole subtraction. At the fixed stationary point,
$(\Pi+H_{\rm ct})Tx_0=0$. The two independent normalized phase directions
provide a sensitive test of the counterterm normalization.

\subsection{Numerical tadpoles and Ward checks}
The bosonic calculation gives
\begin{align}
 (t_\omega,t_\sigma,t_s)/\omega^3
 &=(-0.135193703018,-0.691922627373,-0.008559428597),\\
 (\delta\mu^2,\delta\nu^2,\delta\mu_s^2)/\omega^2
 &=(-0.056330709591,-5.469744089909,-0.034237714387).
\end{align}
For the normalized $\Sigma$ phase, the loop/CT curvatures are
$-2.734872044954$ and $+2.734872044955$; for the $S$ phase they are
$-0.034237714387$ and $+0.034237714387$. Their residuals are below
$1.3\times10^{-13}\omega^2$. Phase and doublet tadpoles vanish to the
checked precision. The full SM-singlet projector equals the projector onto
the three radial and two phase directions.

The raw subtraction $|\delta\nu^2|/|\nu^2|\simeq14.32$ is large.
Its cancellation against a common multiplet curvature is required by the
Ward identity; the ratio is scheme dependent, not by itself a proof of
breakdown. Nevertheless it demands a genuine higher-loop/naturalness
assessment. We do not infer perturbative convergence from a successful
one-loop stationarity test.

\section{Step 2: the complete complex doublet matrix}
\subsection{Canonical complex copies and exact spectral derivatives}
Let $B$ be a $328\times4$ complex orthonormal basis in the complexification
of the real scalar representation with
\begin{equation}
 C_3=0,\quad C_2=3/4,\quad iR_Y=1/2,\quad iR_{T_3}=-1/2.
\end{equation}
Each of the geometric sectors
$(\phi_{\rm hol},\phi_{\rm anti},\Sigma_{\rm hol},\Sigma_{\rm anti})$
has rank one here. Define $Q_R=\sqrt2\Re B$, $Q_I=-\sqrt2\Im B$.
Then $(Q_R,Q_I)$ is an orthonormal real eight-plane and
$e^{-\pi R_Y}Q_R=Q_I$. A Hermitian complex mass matrix $D$ has the real
form
\begin{equation}
 D_{\mathbb R}=\begin{pmatrix}\Re D&-\Im D\\ \Im D&\Re D\end{pmatrix}.
\label{eq:realification}
\end{equation}
We compute all four diagonal, six real mixed and six imaginary mixed
entries. Imaginary jets are also tested independently. The phase of each
column is fixed by its largest canonical coordinate being real positive;
this is a scalar coordinate convention, not an inferred fermionic Clebsch
phase convention.

For an operator $A$ with eigenvalues $\lambda_i$, write
$A_a=\partial_a A$, $A_{ab}=\partial_a\partial_bA$, in its eigenbasis.
For a separated hard branch put $f_i'=F'(\lambda_i)$ on the hard modes
and zero on soft modes. The exact second variation is
\begin{equation}
 \partial_a\partial_b\tr_{\rm hard}F(A)
 =\sum_i f_i'(A_{ab})_{ii}
 +\sum_{ij}f'^{[1]}_{ij}(A_a)_{ij}(A_b)_{ji},
\end{equation}
where $f'^{[1]}_{ij}=(f'_i-f'_j)/(\lambda_i-\lambda_j)$, with the
continuous diagonal/degenerate limit. Hard--soft divided differences
are retained. Dropping them would discard field-dependent eigenvector
mixing. Since $V_0$ is quartic, $H(x)$ is quadratic, so centered first and
second jets are exact polynomials up to floating-point roundoff. Steps
$0.02\omega$ and $0.01\omega$ agree at better than
$10^{-12}\omega^2$ for the independent curvature check.

\subsection{Computed matrix, before and after the finite subtraction}
In units of $\omega^2$, the tree matrix is
\begin{equation}
D_0=\begin{pmatrix}
.016799176&.015202796&0&.000326645\\
.015202796&.013758748&0&0\\
0&0&.285626100&-.010606602\\
.000326645&0&-.010606602&.138405400
\end{pmatrix}.
\end{equation}
The complete hard bosonic correction is
\begin{equation}
\Pi_B=\begin{pmatrix}
-.0194382305&-.0002813614&-.0000003330&.0002014259\\
-.0002813614&-.0213868984&.0000302589&-.0000020155\\
-.0000003330&.0000302589&-2.5883375274&.0000609813\\
.0002014259&-.0000020155&.0000609813&-2.6730027447
\end{pmatrix}.
\end{equation}
The imaginary entries have magnitude below $2.8\times10^{-13}$ for this
real-coupling benchmark; they were computed, not set to zero. Separate
scalar/vector matrices and the full canonical basis are exported in the
JSON ledger. The old tree-light projection
$c_0^\dagger\Pi_Bc_0=-0.0202366510950\omega^2$ is reproduced.

On this four-copy space,
\begin{equation}
D_{\rm ct}=-\frac{\delta\nu^2}{2}\mathrm{diag}(0,0,1,1),\qquad
P_{10}=\mathrm{diag}(1,1,0,0).
\end{equation}
Thus the large common $126$ curvature cannot be omitted while keeping the
fixed VEVs. The retuned bosonic matrix is
\begin{equation}
D_B=D_0+\Pi_B+D_{\rm ct}-\delta\xi_{02}P_{10}.
\label{eq:assembled}
\end{equation}
The $10$ has zero VEV, so this last finite mass tuning does not alter the
three tadpole equations.

\section{Step 3: the corrected bosonic light state}
\subsection{Schur equation and numerical eigenpair}
In the basis $(c_0,Q)$, with $c_0$ the tree zero mode and $Q$ its
orthonormal complement, decompose
\begin{equation}
 D_B=\begin{pmatrix}A&b^\dagger\\b&C\end{pmatrix}.
\end{equation}
For $C>0$, solving the lower block gives the exact algebraic criterion
\begin{equation}
 \boxed{A=b^\dagger C^{-1}b,\qquad
 c_B=\frac{c_0-QC^{-1}b}{\sqrt{1+b^\dagger C^{-2}b}}.}
\label{eq:schur}
\end{equation}
Here ``exact'' refers to the assembled one-loop-truncated matrix, not to
an all-orders field theory. Diagonalizing it resums some higher-order
algebraic terms but supplies none of the missing two-loop diagrams.

The solution has $\delta\xi_{02}/\omega^2=-0.0202627041913$ and
\begin{equation}
 D_B/\omega^2\simeq\begin{pmatrix}
.0176236497&.0149214344&-.0000003330&.0005280705\\
.0149214344&.0126345543&.0000302589&-.0000020155\\
-.0000003330&.0000302589&.4321606176&-.0105456204\\
.0005280705&-.0000020155&-.0105456204&.2002747003
\end{pmatrix}.
\end{equation}
Its eigenvalues and phase-fixed light vector are
\begin{align}
 \operatorname{spec}(D_B)/\omega^2
 &=(-1.1\times10^{-19},\ 0.03025666325,\ 0.19979763559,\ 0.43263922302),\\
 c_B&\simeq(0.6462011668,-0.7631652015,
 0.00001218358,-0.0017108974)^T.
\end{align}
The rotation from $c_0$ is $0.0329814289$ radians,
$\norm{b}=0.000998486637\omega^2$, and both sides of the Schur condition
are $0.00003290941346\omega^2$ with residual
$6.2\times10^{-18}\omega^2$. Projecting the fixed-VEV correction alone
would instead give $\delta\xi_{02}/\omega^2=-0.0202297946947$.

\subsection{A less wasteful, exact heavy-positivity certificate}
The old sufficient norm test gives
\begin{equation}
 \norm{Q^\dagger(D_B-D_0)Q}/\omega^2=0.1465345611
 >\Delta_{\rm tree}/\omega^2=0.0305570303.
\end{equation}
Thus that test fails by a factor $4.79545$; it may not be reported as
passed. Its failure does not prove a negative eigenvalue, because it
compares corrections in all directions with only the smallest mass.
Let $C_0=Q^\dagger D_0Q>0$ and
\begin{equation}
 E=C_0^{-1/2}(C-C_0)C_0^{-1/2}.
\end{equation}
By congruence,
\begin{equation}
 C=C_0^{1/2}(I+E)C_0^{1/2}>0
 \quad\Longleftrightarrow\quad \lambda_{\min}(E)>-1.
\end{equation}
Numerically $\operatorname{spec}(E)\simeq(-0.0109064,0.4393632,0.5234555)$,
so $C\succeq0.9890936C_0>0$. This proves positivity of this heavy block
without treating the $0.1465$ correction as a correction to the lightest
heavy state. The largest relative correction is still about $52\%$;
neither this certificate nor the positive spectrum establishes
perturbative convergence or positivity of the other scalar sectors.

The corrected magnitudes give $|r|=1.18100$, $|s|=0.00602976$ under the
previous scalar-only convention map, while
$w_{126}=|c_3|^2+|c_4|^2=2.92732\times10^{-6}$.
Most of the rotation is within the $10$ pair. A sizable total rotation is
therefore not itself a flavor repair or a growth of the needed $126$
component. Absolute spinor normalization and PS transport must precede a
physical flavor inference.

\section{Fermionic completion in the same tadpole prescription}
For a fixed invertible symmetric $M_R=\sigma f_M$, define $X=M_R^\dagger M_R$.
Here $f_M$ denotes the Majorana-normalized matrix; it is not silently
identified with the unit-bidoublet coupling defined below.
The hard three-Majorana functional at vanishing electroweak background is
\begin{equation}
 V_{1,F}^{\rm hard}=-\frac1{32\pi^2}
 \tr\left[X^2\left(\log\frac{X}{\mu^2}-\frac32\right)\right].
\end{equation}
Since $\partial_\sigma X=2X/\sigma$, spectral functional calculus gives
\begin{equation}
 t_{\sigma,F}=-\frac1{8\pi^2\sigma}
 \tr\left[X^2\left(\log\frac{X}{\mu^2}-1\right)\right],\qquad
 \delta\nu_F^2=\frac{t_{\sigma,F}}{\sigma}.
\label{eq:Ftad}
\end{equation}
In this minimal unbroken-electroweak sector $t_{\omega,F}=t_{s,F}=0$.
This is a partial field derivative; $\mu$ must not vary along with
$\omega$ even though its chosen numerical value is $g_U\omega$.

For doublet-dependent $D(z)=\sum_a z_aY_a$ and the mass matrix
\begin{equation}
 \mathcal M(z)=\begin{pmatrix}0&D(z)\\D(z)^T&M_R\end{pmatrix},
\end{equation}
the previously derived full complex Hessian is
\begin{equation}
 (\Pi_F)_{ab}=-\frac1{8\pi^2}\tr(Y_a^\dagger Y_b W),\qquad
 W=X\left(\log\frac{X}{\mu^2}-1\right).
\end{equation}
Its use in the fixed-VEV prescription requires
\begin{equation}
 \boxed{\Delta D_F=\Pi_F-\frac{\delta\nu_F^2}{2}P_{126}.}
\end{equation}
One must combine this matrix with \eqref{eq:assembled} and solve
\eqref{eq:schur} again. Adding only a fitted scalar number $\eta_Y$ or only
$\Pi_F$ is insufficient. The $Y_a$ and $f_M$ must be derived from the
same two underlying matrices used in the PS fit, with their canonical
Clebsch projections. Synthetic tests verify this functional identity but
do not supply fitted matrices. Three complex-family tests, including a
degenerate Takagi pair and all eight real doublet directions, pass
$11/11$: the radial and Hessian relative errors are below
$2.6\times10^{-10}$ and $8.8\times10^{-8}$, respectively; family/copy
unitary covariance also holds, with $P_{126}$ transformed along with the
copy basis.

\section{Canonical Spin(10) normalization: a second independent repair}
This calculation can be done before a fit. Let $X,Y,Z$ be the Pauli
matrices and set
\begin{equation}
 \Gamma_{2k-1}=Z^{\otimes(k-1)}\otimes X\otimes I^{\otimes(5-k)},\quad
 \Gamma_{2k}=Z^{\otimes(k-1)}\otimes Y\otimes I^{\otimes(5-k)}.
\end{equation}
Then $\{\Gamma_i,\Gamma_j\}=2\delta_{ij}$,
$\Gamma_*=(-i)^5\Gamma_1\cdots\Gamma_{10}$ and
$C=\Gamma_2\Gamma_4\Gamma_6\Gamma_8\Gamma_{10}$.
The tensors $C\Gamma_i$ and $C\Gamma_{i_1\cdots i_5}$ are symmetric,
while the three-form tensors are antisymmetric. Restrict to the chiral
16-dimensional subspace selected by the P1 self-duality. If $U_{Ia}$ is
the actual orthonormal five-form basis, define
\begin{equation}
 (B_{126,a})_{AB}=\sum_{i_1<\cdots<i_5}U_{i_1\cdots i_5,a}
 (C\Gamma_{i_1}\cdots\Gamma_{i_5})_{AB}.
\end{equation}
The ordered-index sum equals the action-card $1/5!$ contraction.
For each of the 45 generators $R$, the verifier checks
\begin{equation}
 R_{16}^TB_a+B_aR_{16}+\sum_b B_b(R_{126})_{ba}=0,
\end{equation}
and the analogous vector identity. Thus the normalization is linked to
the actual scalar tensor, not guessed from another author's notation.

On the unit P1 singlet $\Omega_{126}$ the raw five-form contraction has
rank one with sole singular value $4\sqrt2$. On a unit canonical
doublet, the down-type raw coefficients are $\sqrt2$ for $10$ and
$2/\sqrt3$ for $126$; the signed relative lepton/down Clebsch is $-3$.
Consequently, for the raw action-card matrices $Y_{10},Y_{126}$,
\begin{equation}
 \boxed{h_D=\sqrt2Y_{10},\qquad f_D=\frac2{\sqrt3}Y_{126},\qquad
 f_M=4\sqrt2Y_{126}=2\sqrt6 f_D.}
\end{equation}
These positive CG normalization constants fix absolute magnitudes; the
complex vertices must also carry a common scalar/spinor phase convention.
With unit canonical scalar doublets and the P1 radial $\sigma$,
\begin{equation}
 \boxed{M_R=2\sqrt6\,\sigma f_D=\sigma f_M.}
\end{equation}
Alternatively, if one calls $f_M$ simply $f$, the Dirac $126$ vertex is
$f/(2\sqrt6)$. One cannot use both unit Dirac $f$ and $M_R=\sigma f$
with the same canonical $\sigma$. The action card is corrected accordingly;
the field content and number of independent family matrices do not change.

The selected Clifford chirality has conjugate-matter hypercharges in the
P1 generator orientation. A global conjugation is allowed, but the scalar
holomorphic/antiholomorphic map and fermion states must be conjugated
together. The absolute factors above and the signed relative $-3$ are
checked; they do not yet constitute the phase-resolved projection of all
four scalar copies. Previous external-matrix CW numbers are historical
conditional tests, not physical values after this correction. Since the
heavy spectrum also enters a logarithm, a silent constant rescaling of
the final $\eta_Y$ number is not a valid repair.

The independent $LL$ triplet contraction has raw magnitude $4\sqrt2$,
so the Majorana-normalized convention has
$|M_L|=|f_M\Delta_{\rm can}|$. The actual exported $Y=+1$ triplet
is a $\Sigma$-antiholomorphic direction in the selected conjugate-matter
orientation; its conjugate has unit overlap with the required
holomorphic $LL$ tensor. This verifies the absolute triplet factor as
well. The common complex phase convention with $Y_\nu$ and the light
doublet still must be exported before combining type I and type II.
All $20/20$ Clifford, covariance, normalization and triplet-alignment
checks pass.

\section{Type-II matching requires the full mixed triplet inverse}
The actual tree spectrum has two complex $(1,3,|Y|=1)$ copies from
$54$--$126$ mixing, with
\[
m_{T,1}^2/\omega^2=0.2688312001,\qquad
m_{T,2}^2/\omega^2=0.5416068249.
\]
The $54$ also contains a real $Y=0$ triplet, which is not the direct
$LL$-coupled channel. Let $e_A$ span all neutral real directions of the
$|Y|=1$ triplets, and let $q$ be the unit real neutral light direction
with $x=x_0+v q$ corresponding to $H^0=v/\sqrt2$. Set
\begin{equation}
 K_{AB}=V_0^{(2)}[e_A,e_B],\qquad J_A=V_0^{(3)}[e_A,q,q].
\end{equation}
Then
\begin{equation}
 V=V_0+\frac12z^TKz+\frac{v^2}{2}J^Tz+\cdots,
 \qquad z_{\rm ind}=-\frac{v^2}{2}K^{-1}J+O(v^4/M_T^3).
\end{equation}
Define the canonical fermionic contraction
$\mathcal Y_A^{LL}=\partial M_{LL}/\partial z_A$. In the convention
$m_\nu=v^2C_5/2$,
\begin{equation}
 \boxed{C_5^{II}=-\sum_{AB}\mathcal Y_A^{LL}(K^{-1})_{AB}J_B,
 \qquad C_5^I=-Y_\nu M_R^{-1}Y_\nu^T.}
\label{eq:typeII}
\end{equation}
The fit uses $C_5=C_5^I+C_5^{II}$, unless a controlled bound shows that one
term is negligible. Even $J_{126}=0$ would not suffice to set type II to
zero: $(K^{-1})_{126,54}J_{54}$ can be nonzero. Representation/PQ
allowance of a source is not a numerical computation of its Clebsch.
The finite one-loop seesaw matching must be adapted only after the
additional PS fields have been matched, rather than importing a pure
SMEFT formula at the wrong stage~\cite{typeI,typeII}.

The new scalar-source calculation uses the exported complex $B,c_B$,
with $q=\sqrt2\Re(Bc_B)$, and finds in the phase-fixed complex
$(54,126)$ triplet basis
\begin{equation}
 K_{\mathbb C}/\omega^2=\begin{pmatrix}
 .541600225&-.001341735117i\\ .001341735117i&.268837800
 \end{pmatrix}.
\end{equation}
In real order $(54_R,126_R,54_I,126_I)$,
\begin{align}
 J/\omega&\simeq(0,.002684241669,-.063294492019,0)^T,\\
 \omega z_{\rm ind}/v^2&\simeq(0,-.004700734233,.058421207024,0)^T.
\end{align}
The magnitude of the $126$ response from its direct source is
$0.004992368741/\omega$, while the $54$ source contributes
$0.000291634508/\omega$ with the opposite phase. Suppressing mass mixing
would give $0.004992307015/\omega$: mixing reduces this by about $5.84\%$.
The ten checks include closure of the complete triplet subspace, canonical
real/complex equivalence, two derivative steps, positivity and the
projected field equation. This uses tree $K,V_0'''$ with a
bosonic-improved $q$; it is a mixed-order scalar diagnostic, not a
loop-accurate type-II Wilson coefficient or a fitted neutrino matrix.

\section{Step 4 audit: the Pati--Salam interface must be repaired}
\label{sec:PS}
\subsection{The active census and an exact beta-table discrepancy}
For group order $(4,L,R)$, the declared four active complex scalar parents
have the following indices and quadratic Casimirs:
\begin{center}
\begin{tabular}{lrrr}
\toprule
Parent & $d_{\mathbb C}$ & $(T_4,T_L,T_R)$ & $(C_4,C_L,C_R)$\\
\midrule
$\phi:(1,2,2)$&4&$(0,1,1)$&$(0,3/4,3/4)$\\
$\Sigma:(15,2,2)$&60&$(16,15,15)$&$(4,3/4,3/4)$\\
$\Delta_L:(10,3,1)$&30&$(9,20,0)$&$(9/2,2,0)$\\
$\Delta_R:(\overline{10},1,3)$&30&$(9,0,20)$&$(9/2,0,2)$\\
\bottomrule
\end{tabular}
\end{center}
Using Weyl fermions and complex scalars, the gauge-and-matter terms are
\begin{align}
 a_i={}&-\frac{11}{3}C_i(G)+\frac23\sum_FT_i(F)+\frac13\sum_ST_i(S),\\
 b_{ij}={}&-\frac{34}{3}C_i(G)^2\delta_{ij}
 +\sum_F[2C_j(F)+\tfrac{10}{3}C_i(G)\delta_{ij}]T_i(F)\nonumber\\
 &+\sum_S[4C_j(S)+\tfrac23C_i(G)\delta_{ij}]T_i(S).
\end{align}
For three $16_F$ families, $\sum_F T=(6,6,6)$. Exact rational evaluation
then gives
\begin{equation}
 a_{\rm PS}=(2/3,26/3,26/3),\qquad
 b_{\rm PS}=\begin{pmatrix}
3551/6&249/2&249/2\\1245/2&779/3&48\\1245/2&48&779/3
\end{pmatrix}.
\end{equation}
Yukawa-dependent two-loop gauge terms are additional and cannot be set to
zero in a claimed full fit. The historical table has $a_4=1$ and
$b_{44}=1209/2$. Its difference is exactly one extra complex $(6,1,1)$,
which contributes $\Delta a_4=1/3$, $\Delta b_{44}=38/3$.
The historical coefficients agree with Eq.~(7) of~\cite{BK}, while the
active census follows the two sixes assigned to $M_U$ in its Table~1.
Our independent index calculation exposes this discrepancy; agreement
with a printed table is not independent validation.

\subsection{A matching-scale failure independent of the extra six}
The old two-site code assigns only $\Delta_R$ to the lower scalar
threshold. Both that parent and the lower massive vectors are
$SU(2)_L$ singlets. It therefore computes $\lambda_{I,L}\equiv0$.
But the same code uses the convention
\begin{equation}
 \alpha_{2,\rm SM}^{-1}=\alpha_{L,\rm PS}^{-1}
 -\frac{\lambda_{I,L}}{12\pi},\qquad
 \frac{d\alpha_i^{-1}}{d\log\mu}=-\frac{a_i}{2\pi}+O(\alpha).
\end{equation}
Differentiating demands
\begin{equation}
 \boxed{\lambda'_{I,L}=-6(a_L^{\rm PS}-a_2^{\rm SM})
 =-6(26/3+19/6)=-71,}
\end{equation}
not zero. The residual matching equation changes by
$71/(12\pi)=1.88333349325$ per unit $\log\mu$ at one loop.
This exact mismatch invalidates the physical matching-scale interpretation
of the old root and covariance table. It is not an issue of a few decimal
places and not a proof that the action has no viable point.

\subsection{Same-content logarithmic repair in all six directions}
Retain the four full parents above between the thresholds. In real-scalar
index convention, $I_A=(68,72,72)$ and $I_{\rm full}=(84,84,84)$.
The PS-to-SM map, with SM order $(3,2,1)$, is
\begin{equation}
 \mathcal R(v_4,v_L,v_R)=(v_4,v_L,\tfrac25v_4+\tfrac35v_R),\qquad
 I_h=(0,1,3/5).
\end{equation}
Massive-vector indices follow from adjoint differences:
$I_{V,U}=(4,6,6)$ and $I_{V,I}=(1,0,14/5)$. Subtract each vector site's
Goldstone orbit once from the scalar census. The physical scalar indices
are consequently
\begin{align}
 I_{S,U}&=I_{\rm full}-I_A-I_{V,U}=(12,6,6),\\
 I_{S,I}&=\mathcal R I_A-I_h-I_{V,I}=(67,71,67).
\end{align}
Physical real scalars contribute $+I_S\log(M/\mu)$ to $\lambda$ and the
vector/gauge/ghost sector contributes $-21I_V\log(M/\mu)$. Hence
\begin{align}
 \lambda'_U&=21I_{V,U}-I_{S,U}=(72,120,120),\\
 \lambda'_I&=21I_{V,I}-I_{S,I}=(-46,-71,-41/5).
\end{align}
With $a_{10}=-34/3$, both vector identities hold exactly:
\begin{equation}
 \lambda'_U=-6[a_{10}(1,1,1)-a_{\rm PS}],\qquad
 \lambda'_I=-6[\mathcal R a_{\rm PS}-a_{\rm SM}].
\end{equation}
This closes the six logarithmic counting checks. It does not compute the
finite matching: at nonzero $\omega,\sigma$, mass, parent and Goldstone
projectors need not commute. Substituting the new projector into the old
weighted mass logarithm is not a proof of a consistent two-stage finite
EFT. One must match at a PS-preserving upper background and then in the
lower EFT, or derive an equivalent covariant mixed-operator matching with
the same gauge/tadpole prescription. No replacement $M_I,M_U$ is claimed.

\section{A constrained fit, not an unconstrained matrix optimization}
\subsection{Canonical Yukawa matching and scale covariance}
At either threshold let $Y_a^{(0)}=\mathcal P_a(c,\mu)Y^+$ be the
tree intertwiner/eigenvector projection. Hard full-minus-EFT vertices and
kinetic matrices give, by expanding the inverse square roots of the three
kinetic terms,
\begin{equation}
 Y_a^-=Y_a^{(0)}+\Delta\Gamma_a^{\rm hard}
 -\tfrac12\left[(\Delta Z_L)^T Y_a^{(0)}+Y_a^{(0)}\Delta Z_R
 +\sum_b(\Delta Z_\phi)_{ba}Y_b^{(0)}\right]+O(\hbar^2).
\end{equation}
The corresponding logarithmic check is
\begin{equation}
 \frac{d\Delta Y}{d\log\mu}=\beta_Y^-
 -\mathcal P\beta_Y^+-\dot{\mathcal P}Y^++O(\hbar^2).
\end{equation}
The derivative of the doublet projection matters. The scalar number
$\eta_Y$ is not a replacement for vertex or fermion wave-function
matching. Standard PS equations~\cite{MOR} are useful regression targets,
but their minimal example retains $\Delta_R$ without $\Delta_L$ and has a
different lower Higgs EFT. The full P54 left/right matrix beta functions
and absolute canonical intertwiners must be supplied explicitly.

At nondegenerate type-I thresholds write
$M_R=U^*\mathrm{diag}(M_i)U^\dagger$ and $y_i=(Y_\nu U)_{:i}$. In the
convention $m_\nu^I=-v^2\kappa/2$, sequential removal gives
$\kappa^- =\kappa^++y_iy_i^T/M_i$ at $\mu_i\simeq M_i$. Their order must
follow the fitted Takagi masses. A common threshold $M_I$ for all three is
an extra approximation, not a convention.

\subsection{Eliminate nuisance eigenvectors rather than fitting them freely}
Let $\theta$ contain invariant scalar couplings, VEVs, gauge parameters and
the two family matrices modulo common unitary family rotations. In a
nonsingular Takagi patch one may take $h$ diagonal nonnegative and retain
complex symmetric $f$; degenerate patches require their residual quotient.
Define
\begin{equation}
 \mathcal D(\theta,\delta\xi)=D_0+\Pi_B+H_{{\rm ct},B}
 +\Pi_F[h,f]+H_{{\rm ct},F}[f]-\delta\xi P_{10}.
\end{equation}
At each candidate, solve the tadpole equations and the unique light-state
constraint, rather than allowing $c$ and $\eta_Y$ to move independently of
the action. For a simple normalized eigenpair,
\begin{equation}
 d\lambda=c^\dagger(d\mathcal D)c,\qquad
 Q_c\,dc=-[Q_c(\mathcal D-\lambda)Q_c]^{-1}
 Q_c(d\mathcal D)c,\quad c^\dagger dc=0.
\end{equation}
These are obtained by differentiating $\mathcal D c=\lambda c$ and
projecting onto $c$ and its complement. They supply implicit derivatives
for a constrained optimizer; eigenvectors need not be refitted as free
spurions. In particular,
$\partial\lambda/\partial\delta\xi=-c^\dagger P_{10}c<0$ on this branch,
so the light tuning is locally well posed as long as the gap remains open.

The physical likelihood is then
\begin{equation}
 \chi^2_{\rm prof}(\vartheta)=
 \inf_{\theta\in\mathcal C(\vartheta)}
 [\mathcal O_{\rm EFT}(\theta)-d]^T\mathsf C^{-1}
 [\mathcal O_{\rm EFT}(\theta)-d]+\chi^2_{\rm nuisance},
\end{equation}
where $\mathcal C$ imposes same-action stationarity, exactly one light
doublet, heavy positivity, a predeclared perturbativity/control criterion,
consistent finite gauge/Yukawa matching, and both terms of
\eqref{eq:typeII}. The quantity $\eta_Y=c^\dagger\Pi_Fc/\omega^2$ is a
derived correlated nuisance once actual matrices are fitted, not an
arbitrary independent knob. Experimental covariance and any threshold
priors must be stated. An undefined EFT map is not an empty feasible
domain; the optimizer must report ``not ready'', not a spurious best fit
or a proof of exclusion.

\section{Reproducibility and precise completion boundary}
The full-doublet implementation is
\path{route_f/code/verify_p54_full_doublet_cw.py}, with JSON/Markdown
outputs of the same stem. It exports the full canonical basis, separate
scalar/vector matrices, radial counterterms, phase Ward residuals,
retuned eigenpair and relative positivity certificate. Its 18 checks pass.
The exact PS census is implemented in
\path{route_f/code/verify_p54_ps_active_census.py}; its 14 audit tests pass
while the two historical physical consistency gates fail. Historical P2
outputs are retained for reproducibility and explicitly superseded in the
roadmap, not overwritten to conceal the discrepancy.

Further independent verifiers are
\path{route_f/code/verify_p54_spinor_intertwiners.py} ($20/20$),
\path{route_f/code/verify_p54_fermion_tadpole.py} ($11/11$), and
\path{route_f/code/verify_p54_typeii_triplet_source.py} ($10/10$).
The numerical implementation uses Python 3.9, NumPy 2.0.2, SciPy 1.13.1
and JAX/JAXLIB 0.4.30 in double precision. Expensive Hessians are cached
under \path{tmp/p54_full_doublet_cw}, keyed by the P1 action source,
parameters and exact field coordinates; an unchanged replay uses the
same tensors instead of repeating a scan. Source hashes are exported.

Completed here: the invariant tadpole rule; numerical bosonic radial/phase
tests; the complete bosonic complex doublet matrix; its corrected
one-loop-truncated eigenpair; the exact six-direction PS logarithmic repair.
The fermionic rule is a functional identity awaiting the same fitted
matrices and canonical contractions. Still open: finite staged matching,
complete P54 matrix RGEs, total quantum scalar stability/loop control,
the self-consistent flavor/seesaw fit and physical pole observables.
The independent soliton/GTA line does not block this work.

\begin{thebibliography}{9}
\bibitem{BK} K.~S.~Babu and S.~Khan,
\emph{A Minimal Non-Supersymmetric SO(10) Model: Gauge Coupling
Unification, Proton Decay and Fermion masses},
\href{https://arxiv.org/abs/1507.06712}{arXiv:1507.06712}.
\bibitem{MOR} D.~Meloni, T.~Ohlsson and S.~Riad,
\emph{Renormalization Group Running of Fermion Observables in an
Extended Non-Supersymmetric SO(10) Model},
\href{https://arxiv.org/abs/1612.07973}{arXiv:1612.07973}.
\bibitem{typeI} D.~Zhang and S.~Zhou,
\emph{Complete One-loop Matching of the Type-I Seesaw Model onto the
Standard Model Effective Field Theory},
\href{https://arxiv.org/abs/2107.12133}{arXiv:2107.12133}.
\bibitem{typeII} X.~Li, D.~Zhang and S.~Zhou,
\emph{One-loop Matching of the Type-II Seesaw Model onto the
Standard Model Effective Field Theory},
\href{https://arxiv.org/abs/2201.05082}{arXiv:2201.05082}.
\end{thebibliography}
\end{document}
